\begin{abstract}
A medical AI system misclassifies a critical diagnosis due to a single typo—a failure pattern repeated millions of times daily as transformer models struggle with noisy real-world data. Through forensic layer-wise analysis of five transformer architectures processing 300,000 samples, we uncover why: models undergo critical vulnerability transitions at layers 3 and 8, marking boundaries between surface, syntactic, and semantic processing phases. These "fault lines" determine whether models catastrophically fail or remarkably recover. Our investigation reveals that RoBERTa achieves near-perfect robustness (0.988) while ELECTRA collapses (0.527), that character-level noise shows 85\% recovery while syntax perturbations cause 78\% degradation, and that vulnerability patterns transfer across architectures with 61.1\% correlation. Most remarkably, by exploiting these phase boundaries through strategic layer dropout, we achieve 3.1× inference speedup and 60\% cost reduction while maintaining 95\% performance. These findings transform our understanding of transformer fragility, providing immediate deployment guidelines for noise-critical applications and inspiring phase-aware architectures that turn vulnerability into efficiency—essential steps toward AI systems worthy of the trust we place in them.
\end{abstract}